
\needspace{4\baselineskip}
\color{uniblue}{\subsection{Комбинированный пусковой орган напряжения (КПОН)}}
\color{black}

\begin{enumerate}[label=\arabic{section}.\arabic{subsection}.\arabic*, labelsep=4pt, leftmargin=0pt, itemindent=57pt]

\item
Функция КПОН формирует разрешающий сигнал для срабатывания внешних ступеней МТЗ, которые в нормальном режиме заблокированы или работают с загрубленной уставкой. В устройстве комбинированный пусковой орган напряжения представлен функцией «КПОН».

\item
В качестве воздействующих величин ИО КПОН использует действующие значения линейных напряжений и напряжение ОП.

\item
Функция содержит в себе ИО минимального напряжения, реагирующий на снижение максимального значения из подведенных линейных напряжений ниже уставки, а также ИО максимального напряжения, реагирующий на превышение напряжения обратной последовательности уставки.

\item
Посредством положения уставки SGF2 «Реж\_пуска» осуществляется выбор алгоритма работы функции:
\begin{itemize}
    \item если уставка в положении «По Uмин» - алгоритм работы функции построен по снижению минимального из линейных напряжений ниже уставки;
    \item если уставка в положении «По Uмин и U2макс» - алгоритм работы функции комбинированный и построен по снижению минимального из линейных напряжений ниже уставки или при повышении напряжения ОП уставки.
\end{itemize}

\item
При помощи уставки SGF3 «Реж\_БНН», контролируется необходимость блокировки КПОН при срабатывании БНН.

\item
Ввод функции в работу на этапе параметрирования устройства осуществляется программным переключателем SGF1 «Ввод\_функции». Информация о введенной функции отображается выходным сигналом «КПОН: Ввод».

\item
Функция может быть выведена из работы оперативно путем активации сигнала «ОВ~КПОН», что характеризуется активным состоянием сигнала «КПОН: Оперативный вывод» на выходе функции. 

\item
Алгоритм работы функции выполнен в соответствии с рисунком \ref{kpon:img1}. В таблице \ref{kpon:tbl1} приведены параметры, необходимые для настройки функции.

\vspace{3mm}
\begin{figure}[H]
\centering
\includegraphics[width=0.9\textwidth,height=0.9\textheight,keepaspectratio]{\fbpath/212. КПОН ТН/_latex/img/kpon.pdf}
\captionof{figure}{Функциональная схема логики КПОН}\label{kpon:img1}
\end{figure}

\end{enumerate}
